% !TEX root = ../algebra_02.tex

\section{Порядки элементов и циклические подгруппы}

\begin{Definition}{Порядок элемента}{}
	Если существует натуральное число $n \in \NN$ такое, что $g^n = e$, то говорят, что элемент $g$ имеет конечный порядок. Порядком элемента $g$ называется минимальное такое число:
	\[
	\operatorname{ord} g = \min\{n \in \NN \mid g^n = e\}
	\].
	Если такого $n$ не существует, то $\operatorname{ord} g = \infty$.
\end{Definition}

\begin{Proposition}{}{}
	Пусть $\operatorname{ord} g = d < \infty$. Тогда $\{m \in \ZZ \mid g^m = e\} = d\ZZ.$
	То есть $g^m = e \iff d \mid m.$
\end{Proposition}

\begin{proof}
Разделим $m$ на $d$ с остатком: $m = dq + r, \qquad 0 \le r < d.$
Тогда $g^m = e \implies g^{dq+r} = e \implies (g^d)^qg^r = e.$
Так как $g^d = e$, получаем $g^r = e$. Поскольку $d$ — минимальная натуральная степень, дающая $e$, а $0 \le r < d$, единственный возможный вариант — $r = 0$. Следовательно, $m = dq$, то есть $d \mid m$.
\end{proof}

\begin{Definition}{Циклическая подгруппа}{}
	Множество $\langle g \rangle = \{g^n \mid n \in \ZZ\}$ называется циклической подгруппой, порождённой элементом $g$.
\end{Definition}

\begin{Definition}{Циклическая группа}{}
	Если существует такой элемент $g \in G$, что $\langle g \rangle = G,$
	то группа $G$ называется циклической, а элемент $g$ — её образующим, или порождающим, элементом.
\end{Definition}

Размер (порядок) циклической подгруппы напрямую зависит от порядка порождающего элемента:
\begin{Proposition}{}{}
	Пусть $g \in G$.
	\begin{enumerate}
		\item Если $\operatorname{ord} g = \infty$, то $|\langle g \rangle| = \infty$ и все степени $g^i$ различны.
		\item Если $\operatorname{ord} g = n < \infty$, то $|\langle g \rangle| = n$ и $\langle g \rangle = \{g^0,g^1,\dots,g^{n-1}\}$.
	\end{enumerate}
\end{Proposition}

\begin{proof}
	Если $\operatorname{ord} g = \infty$, предположим, что $g^i = g^j$ при $i > j$. Тогда $g^{i-j} = e$, и так как $i-j > 0$, это противоречит бесконечности порядка. Значит, все степени различны.

	Если $\operatorname{ord} g = n < \infty$, для любого $m \in \ZZ$ разделим $m$ на $n$ с остатком: $m = nq + r$, где $0 \le r < n$. Тогда $g^m = (g^n)^q g^r = e^q g^r = g^r$, поэтому множество степеней ограничивается $\{g^0,\dots,g^{n-1}\}$. Все они различны, так как иначе существовала бы меньшая положительная степень $i-j < n$, дающая единицу, что противоречило бы минимальности порядка $n$.
\end{proof}

\begin{Proposition}{}{}
	Пусть $|G| = n < \infty$. Тогда $G$ является циклической тогда и только тогда, когда в ней существует элемент порядка $n$.
\end{Proposition}

\subsubsection*{Доказательство}

\begin{itemize}
	\item[$\Rightarrow$)] Если $G$ циклическая, то по определению $G = \langle a \rangle$ для некоторого $a \in G$. Тогда $\operatorname{ord} a = |\langle a \rangle| = |G| = n.$
	
	\item[$\Leftarrow$)] Пусть существует $g \in G$ такой, что $\operatorname{ord} g = n$. Тогда $|\langle g \rangle| = n.$
	Так как $\langle g \rangle \subseteq G$ и их порядки совпадают, то $\langle g \rangle = G.$
	Следовательно, $G$ циклическая.
\end{itemize}